% Einheit 1 von 5 – Winkel messen und zeichnen (bank/winkel-dreiecke/e1.jsonl, zusammenbau v0.1)
%% TODO Zweigzeile Teil 1: was hier gelernt wird (Satz in Schülersprache aus dem Einheitstitel) – Einheit 1
\einheitenkopf[e1]{Einheit 1 von 5 · Winkel messen und zeichnen}
\zweigzeile{ab Kl. 5 · P10}
\setcounter{aufgabe}{7}

%% TODO Titel: Ich-kann-Satz fehlt in der Bank (Platzhalter: Kettenname) – Nr. 8
%% TODO Anweisung: ein Satz über der ersten Teilaufgabe fehlt in der Bank – Nr. 8
\begin{aufgabe}{Welche Skala?}
\begin{teile}
\teil Nicht ablesen: Gilt beim Messen von α am Geodreieck die kleinere oder die größere der beiden Zahlen? \\ \kreuz{die kleinere Zahl} \\ \kreuz{die größere Zahl}

\winkel[5]{38}{\alpha}
\end{teile}
\end{aufgabe}

%% TODO Titel: Ich-kann-Satz fehlt in der Bank (Platzhalter: Kettenname) – Nr. 9
%% TODO Anweisung: ein Satz über der ersten Teilaufgabe fehlt in der Bank – Nr. 9
\begin{aufgabe}{Winkel messen}
%% TODO \rechenplatz{n}: Schrittzahl des Grundfalls fehlt in der Bank (nur bei fester Schreibform, 2.3 b)
\begin{teile}
\teil Welche Winkelart hat α? Nicht messen. \\ \kreuz{spitz} \\ \kreuz{recht} \\ \kreuz{stumpf} \\ \kreuz{gestreckt} \\ \kreuz{überstumpf}

\winkel[4]{27}{\alpha}
\teil Miss den Winkel α mit dem Geodreieck.

\winkel[5]{37}{\alpha}
α = \leerfeld °
\teil Miss den stumpfen Winkel α mit dem Geodreieck.

\winkel[5]{118}{\alpha}
α = \leerfeld °
\teil Zeichne an den Strahl mit dem Scheitel S einen Winkel von $55^\circ$.

\winkelstrahl
\teil Bestimme den überstumpfen Winkel α. Miss dazu den Rest bis zum Vollwinkel.

\winkel[4]{215}{\alpha}
α = \leerfeld °
\teil Schreibe den Winkel α im Dreieck ABC mit drei Buchstaben ($\angle$…). Welcher Punkt ist der Scheitel?

\dreieck{(0,0)}{(5,0)}{(1.5,3)}{a}{b}{c}{\alpha}{\beta}{\gamma}
\teil Die Winkel $34^\circ$ und $47^\circ$ liegen am selben Scheitel nebeneinander. Wie groß ist der Winkel, den beide zusammen bilden? \leerfeld °
\teil Ein Winkel von $96^\circ$ ist durch einen Strahl in zwei Teilwinkel geteilt. Der eine misst $38^\circ$. Wie groß ist der andere? \leerfeld °
\teil Ein Wanderweg ist 1,2 km lang. Nach 450 m macht Ella eine Pause. Wie viele Meter hat sie noch vor sich? \leerfeld[m]
\teil Vom Punkt A aus sieht Jana die Spitze T eines Turms unter einem Winkel von $41^\circ$ zur waagerechten Grundlinie, die Spitze F der Fahne auf dem Turm unter $44^\circ$. Wie groß ist der Winkel γ zwischen den Sichtlinien AT und AF? \hfill (P10 2018 OS) \\ γ = \leerfeld °
\end{teile}
\end{aufgabe}

%% TODO Titel: Ich-kann-Satz fehlt in der Bank (Platzhalter: Kettenname) – Nr. 10
%% TODO Anweisung: ein Satz über der ersten Teilaufgabe fehlt in der Bank – Nr. 10
\begin{aufgabe}{Winkel schätzen}
\begin{teile}
\teil Schätze, ohne zu messen: Wie groß ist α im Vergleich zu einem rechten Winkel? \\ \kreuz{ein Viertel} \\ \kreuz{die Hälfte} \\ \kreuz{drei Viertel}

\winkel[5]{23}{\alpha}
\end{teile}
\end{aufgabe}

%% TODO Titel: Ich-kann-Satz fehlt in der Bank (Platzhalter: Kettenname) – Nr. 11
%% TODO Anweisung: ein Satz über der ersten Teilaufgabe fehlt in der Bank – Nr. 11
\begin{aufgabe}{Winkel messen – Fehler finden, Begründen, Anwendung}
\begin{teile}
\teil Tim hat den Winkel α gemessen und notiert: $\alpha = 144^\circ$. Finde den Fehler und gib den richtigen Wert an.

\winkel[5]{36}{\alpha}
\teil Warum heißt ein Winkel von $91^\circ$ stumpf, ein Winkel von $89^\circ$ aber spitz? Begründe.
\teil Eine Parkhausschranke dreht sich beim Öffnen aus der waagerechten Lage, bis sie senkrecht steht. Nach zwei Sekunden hat sie sich um $54^\circ$ gedreht. Um wie viel Grad muss sie sich noch drehen? \leerfeld °
\end{teile}
\end{aufgabe}
